\chapter{The Structural Verification Layer}\label{ch:verification}

\section{Motivation: Verifying Before Emitting, Not After}\label{sec:verify-motivation}

\section{Claim Decomposition of the Draft Answer}\label{sec:claim-decomposition}

\section{Two-Stage Claim Checking}\label{sec:two-stage}
% Call these the STRUCTURAL CHECK and the ENTAILMENT CHECK. Never "Tier 1 /
% Tier 2" --- "tier" is reserved for the groundedness metric in
% Chapter~\ref{ch:setup}.
\subsection{The Structural Check}\label{sec:structural-check}
\subsection{The Entailment Check}\label{sec:entailment-check}
\subsection{Why the Structural Check Runs First}

\section{Repair Policies for Unsupported Claims}\label{sec:repair}
\subsection{Targeted Re-Exploration Under Remaining Budget}
\subsection{Answer Rewriting and Hedging}
\subsection{Iteration Cap and the Draft-Only Fallback}

\section{Entity Filtering of the Final Answer}\label{sec:entity-filtering}

\section{Output Contract: Answer Paired with Supporting Triples}\label{sec:output-contract}

\section{A Worked Example}\label{sec:verify-example}

\section{Failure Modes by Construction: Wrongful Rejection and Wrongful
  Acceptance}\label{sec:verify-failure-modes}
% Analytical, derived from the design rather than the data --- but
% Section~\ref{sec:verifier-errors} has empirical specimens for both polarities,
% so name the mechanism each one instantiates and forward-reference them. Use
% the polarity vocabulary fixed in Section~\ref{sec:polarities}; no fn/fp here.

\endinput
